\documentclass[11pt]{article}
\usepackage[margin=1in]{geometry}
\usepackage{amsmath,amssymb,amsthm,bm}
\usepackage{microtype}
\usepackage{booktabs}
\usepackage[hidelinks]{hyperref}
\usepackage[backend=biber,style=numeric,sorting=none,natbib=true]{biblatex}
\addbibresource{references.bib}
\usepackage{enumitem}
\usepackage{xcolor}

\newtheorem{proposition}{Proposition}
\newtheorem{remark}{Remark}
\newcommand{\dd}{\mathrm{d}}
\newcommand{\mpl}{M_{\rm Pl}}
\newcommand{\Geff}{\mathcal G}
\newcommand{\Ddark}{\mathcal D}
\newcommand{\Pontryagin}{\mathcal P}
\newcommand{\ZE}{\mathbb Z_p}

\title{Finite-State Entropic Defects in Curved Spacetime:\\
Curvature, Rotation, and Dark-Sector Triggering}
\author{Working research draft}
\date{2 October 2026}

\begin{document}
\maketitle

\begin{abstract}
We formulate a compact effective field theory motivated by a ``universe weather'' picture but expressed only through covariant quantities.  The new sector is a hidden Abelian gauge theory containing a charge-$p$ Higgs field.  Curvature, rotational curvature, and the dark-sector stress-tensor trace shift the Higgs effective mass.  When the corresponding phase exists, $U(1)$ is Higgsed to a residual local $\mathbb Z_p$ gauge symmetry supporting quantized flux defects.  Periodic arithmetic is therefore retained only through the exact holonomy $\exp(2\pi i n/p)$, while the proposed bulk pressure is derived from the hidden-sector stress tensor.  We distinguish a local negative-mass region from a genuine instability: for Schwarzschild we derive the exact radial spectral problem and a Rayleigh--Ritz sufficient criterion for tachyonic onset.  For Kerr we add a parity-even rotational trigger proportional to the square of the Pontryagin density.  Using the exact Kerr invariants, we derive the slow-spin deformation of the local activation surface,
\[
\frac{r_{\rm w}(\theta)}{r_0}=1+\frac{a^2\cos^2\theta}{r_0^2}\left(-\frac72+6\eta_\Omega\right)+O(a^4/r_0^4),
\qquad
\eta_\Omega=\frac{\alpha_\Omega m_0^2 M_G^4}{\alpha_G^2M_\Omega^6},
\]
so the leading polar deformation changes sign at the mass-independent threshold $\eta_\Omega=7/12$.  We also retain the analytic dark-halo eligibility condition for an NFW profile.  Existing Gauss--Bonnet and Chern--Simons scalarization, environment-dependent defects, and discrete gauge theory make clear that none of these ingredients individually is new.  The novelty hypothesis is the combined finite-state gauge phase with curvature-, rotation-, and dark-environment-dependent activation together with a $p$-dependent topological observable.  The construction remains an EFT research program, not a microscopic theory of quantum gravity.
\end{abstract}

\section{Scope and guiding constraints}

The original motivation is an analogy between gravitational systems and weather patterns: flow, pressure, vorticity, calm cores, and violent surrounding structure.  An analogy is not a theory.  A publishable construction must be coordinate invariant, possess a controlled action, recover known gravity in its domain of validity, and generate observables not introduced by definition.

Three constraints therefore guide the model.
\begin{enumerate}[leftmargin=1.7em]
    \item The ``wind'' is not promoted to a preferred physical vector.  Flow decompositions of a spacetime congruence are useful kinematics, but a Painlev\'e--Gullstrand river velocity is coordinate dependent.
    \item ``Entropic pressure'' is not postulated as a scalar force.  It is defined operationally from the stress tensor of a new sector.  A thermodynamic entropy interpretation is a separate hypothesis to be earned.
    \item Periodic arithmetic enters only where it produces an invariant observable.  We use a residual $\mathbb Z_p$ gauge structure.  We do \emph{not} use finite-field multiplication, because no physical role for it is currently derived.
\end{enumerate}

The construction should be regarded as an effective field theory (EFT), not as a claim that spacetime is literally a fluid.


\subsection{What survives from the weather analogy}
A single circular orbit is not a spacetime vortex: Schwarzschild spacetime is nonrotating and nevertheless admits circular geodesics.  The analogy becomes mathematically useful only after distinguishing matter circulation from rotational curvature.  For a Newtonian Keplerian continuum,
\begin{equation}
v_\phi(r)=\sqrt{\frac{GM}{r}},\qquad
\Omega=\sqrt{\frac{GM}{r^3}},
\end{equation}
and the vertical vorticity is
\begin{equation}
\omega_z=\frac1r\frac{\dd}{\dd r}(rv_\phi)=\frac12\Omega.
\end{equation}
This is genuine vortex-like \emph{matter} kinematics, as in a disk, not a new law of gravity.  For spacetime itself we instead use curvature invariants.  In particular, the parity-odd Pontryagin density vanishes for Schwarzschild but is generally nonzero for Kerr, providing a covariant rotational-curvature diagnostic.  This separation prevents the motivating ``orbit as tornado'' picture from being promoted into a false identity.

\section{Why finite-field-valued pressure is rejected}

A finite field $\mathbb F_p$ cannot be an ordered field.  This elementary obstruction matters because pressure requires an ordered real quantity.

\begin{proposition}
No finite field admits an order compatible with field addition and multiplication.
\end{proposition}

\begin{proof}
Assume $1>0$.  Closure of an ordered field under addition implies
\[
\underbrace{1+\cdots+1}_{p}>0.
\]
In $\mathbb F_p$ the left-hand side is $0$, giving $0>0$, a contradiction.
\end{proof}

Thus $P_E\in\mathbb F_p$ is discarded.  The finite structure instead labels internal states or gauge holonomies, while energy density and pressure remain real-valued.

\section{Minimal covariant model}

We work in four dimensions with metric signature $(-,+,+,+)$ and natural units $\hbar=c=1$.  The action is
\begin{equation}
S=S_{\rm EH}+S_E+S_D+S_{\rm vis},
\end{equation}
with
\begin{equation}
S_{\rm EH}=\int \dd^4x\sqrt{-g}\,\frac{\mpl^2}{2}R,
\end{equation}
and an ``entropy-sector'' Abelian Higgs action
\begin{equation}
S_E=\int \dd^4x\sqrt{-g}\left[
-\frac14 F_{\mu\nu}F^{\mu\nu}
-(D_\mu H)^\ast D^\mu H
-V_{\rm eff}(H;x)
\right].
\label{eq:SE}
\end{equation}
Here
\begin{equation}
D_\mu H=(\nabla_\mu-i p g_E A_\mu)H,
\qquad p\in\mathbb Z,\quad p\ge2,
\end{equation}
and
\begin{equation}
V_{\rm eff}=m_E^2(x)|H|^2+\lambda |H|^4,
\qquad \lambda>0.
\label{eq:Veff}
\end{equation}
The invariant environmental mass is
\begin{equation}
\boxed{
 m_E^2(x)=m_0^2
 -\frac{\alpha_G}{M_G^2}\,\Geff(x)
 -\frac{\alpha_\Omega}{M_\Omega^6}\,\Pontryagin(x)^2
 -\frac{\alpha_D}{M_D^2}\,\Ddark(x)
}
\label{eq:meff}
\end{equation}
with
\begin{equation}
\Geff=R_{\mu\nu\rho\sigma}R^{\mu\nu\rho\sigma}
-4R_{\mu\nu}R^{\mu\nu}+R^2,
\end{equation}
\begin{equation}
\Pontryagin=\frac12\epsilon^{\alpha\beta\gamma\delta}
R^\mu{}_{\nu\gamma\delta}R^\nu{}_{\mu\alpha\beta},
\end{equation}
and
\begin{equation}
\Ddark\equiv- T^{(D)\mu}{}_{\mu}.
\end{equation}
The $\Pontryagin^2$ term is parity even and vanishes in Schwarzschild.  Since $[\Geff]=[\Pontryagin]=M^4$, both $\Geff/M_G^2$ and $\Pontryagin^2/M_\Omega^6$ have mass dimension two.  The latter is a higher-order EFT operator and must remain perturbatively controlled; it is included because it distinguishes rotational curvature without introducing a preferred-frame vector.
For pressureless dark matter, $\Ddark\simeq\rho_D>0$.

Equation~\eqref{eq:meff} contains three logically distinct hypotheses.  The Gauss--Bonnet term is a curvature-triggered EFT interaction; the $\Pontryagin^2$ term is a parity-even rotational-curvature trigger; and the dark term is a \emph{dark-selective} interaction.  If the last were universal over visible matter it would face strong fifth-force and equivalence-principle constraints.  A microscopic completion could generate the latter through a dark-sector conformal metric $\tilde g^{(D)}_{\mu\nu}=A^2(|H|)g_{\mu\nu}$, whose small-field expansion produces a coupling to the dark stress-tensor trace.

\subsection{Environmental Higgsing}

Let $s=|H|^2$.  Then
\begin{equation}
V_{\rm eff}(s)=m_E^2 s+\lambda s^2,
\end{equation}
and
\begin{equation}
\frac{\partial V_{\rm eff}}{\partial s}=m_E^2+2\lambda s.
\end{equation}
Hence
\begin{equation}
|H|^2_{\rm vac}=
\begin{cases}
0, & m_E^2\ge0,\\[1mm]
-\dfrac{m_E^2}{2\lambda}, & m_E^2<0.
\end{cases}
\label{eq:vev}
\end{equation}
When $m_E^2<0$, a field of gauge charge $p$ condenses and leaves a local residual $\mathbb Z_p$ gauge symmetry.  This is the precise location of the periodic arithmetic.

The activation criterion is therefore
\begin{equation}
\boxed{
\frac{\alpha_G}{M_G^2}\Geff
+\frac{\alpha_\Omega}{M_\Omega^6}\Pontryagin^2
+\frac{\alpha_D}{M_D^2}\Ddark
>m_0^2.
}
\label{eq:activation}
\end{equation}
This is not yet a defect-formation theorem, nor even a proof of linear instability.  It identifies where the local potential favors the broken phase.  Gradient energy, horizon boundary conditions, and the global spectrum can prevent condensation even when $m_E^2<0$ locally; actual defect nucleation additionally depends on nonlinear history and boundary conditions.

\section{Finite arithmetic as an observable: flux and linking phase}

Consider a straight finite-energy vortex with winding $n\in\mathbb Z$.  Far from its core,
\begin{equation}
H\longrightarrow \frac{v}{\sqrt2}e^{in\varphi},
\end{equation}
where $v^2=-m_E^2/\lambda$ in an approximately homogeneous broken region.  Finite energy requires $D_\mu H\to0$ around a large loop $C$.  Therefore
\begin{equation}
 p g_E\oint_C A_\mu\dd x^\mu=2\pi n,
\end{equation}
so the hidden magnetic flux is
\begin{equation}
\boxed{
\Phi_E(n)=\oint_C A_\mu\dd x^\mu
=\frac{2\pi n}{p g_E}.
}
\label{eq:flux}
\end{equation}
A unit-charge hidden probe acquires the Aharonov--Bohm phase
\begin{equation}
\boxed{
\exp\left(i g_E\Phi_E(n)\right)
=\exp\left(\frac{2\pi i n}{p}\right).
}
\label{eq:AB}
\end{equation}
Consequently
\begin{equation}
\mathcal A_{n+p}=\mathcal A_n
\end{equation}
for observables that depend only on the discrete holonomy.  This is an exact, non-decorative appearance of periodic arithmetic.

At energies below the Higgs and gauge masses, the topological sector is represented by a $\mathbb Z_p$ BF theory after conventional field rescalings,
\begin{equation}
S_{\rm BF}=\frac{p}{2\pi}\int B\wedge \dd A,
\label{eq:BF}
\end{equation}
which encodes the mutual phase between line and surface operators.  This low-energy representation is standard; it is not itself a novelty claim \citep{KapustinSeiberg2014,GaiottoEtAl2015}.

\section{Entropic pressure is derived from stress, not imposed}

The stress tensor of Eq.~\eqref{eq:SE} is
\begin{align}
T^{(E)}_{\mu\nu}={}&F_{\mu\alpha}F_\nu{}^\alpha
-\frac14 g_{\mu\nu}F_{\alpha\beta}F^{\alpha\beta}
+(D_\mu H)^\ast D_\nu H+(D_\nu H)^\ast D_\mu H\nonumber\\
&-g_{\mu\nu}\left[(D_\alpha H)^\ast D^\alpha H+V_{\rm eff}\right].
\end{align}
For an observer with four-velocity $u^\mu$, define
\begin{equation}
h_{\mu\nu}=g_{\mu\nu}+u_\mu u_\nu,
\end{equation}
\begin{equation}
\boxed{
\rho_E=T^{(E)}_{\mu\nu}u^\mu u^\nu,
\qquad
P_E=\frac13 h^{\mu\nu}T^{(E)}_{\mu\nu}.
}
\label{eq:pressure}
\end{equation}
Equation~\eqref{eq:pressure} is the only quantity in this paper called bulk entropic pressure.  It is real, local, observer-dependent in the standard relativistic fluid sense, and calculable once a solution is specified.

The original ``low-pressure eye'' intuition becomes a test rather than an axiom:
\begin{equation}
P_E(r=0)<P_E(r\gg \xi_E)?
\label{eq:pressuretest}
\end{equation}
For an Abelian-Higgs vortex this inequality is not automatic because scalar potential, gradient, and gauge-field stresses compete.  A numerical vortex solution is therefore a necessary next calculation.

\section{Schwarzschild eligibility and the true instability problem}

Curvature-induced scalarization is already established in other scalar theories \citep{DonevaYazadjiev2018}; the calculation here only establishes the threshold scaling of the present model.

For Schwarzschild in geometrized units $G=c=1$,
\begin{equation}
\Geff=\frac{48M^2}{r^6}.
\end{equation}
Ignoring the dark term, the local eligibility surface $m_E^2=0$ obeys
\begin{equation}
m_0^2=\frac{48\alpha_G M^2}{M_G^2 r_G^6}.
\end{equation}
Therefore
\begin{equation}
\boxed{
r_G=
\left(
\frac{48\alpha_G M^2}{m_0^2M_G^2}
\right)^{1/6}.
}
\label{eq:rG}
\end{equation}
The scaling
\begin{equation}
r_G\propto M^{1/3}
\end{equation}
is a direct consequence of a quadratic curvature invariant.

For activation outside the Schwarzschild horizon $r_H=2M$, one needs
\begin{equation}
r_G>2M,
\end{equation}
which is equivalent to
\begin{equation}
\boxed{
m_0^2M_G^2M^4<\frac{3}{4}\alpha_G.}
\label{eq:outside}
\end{equation}
This inequality is a useful kill criterion: for a proposed parameter set it immediately determines whether the curvature-triggered phase is externally accessible or hidden behind the horizon.

The EFT expansion also requires that neglected higher-curvature operators remain suppressed.  A sufficient schematic condition is
\begin{equation}
\frac{|\Geff|}{M_G^4}\ll1
\end{equation}
at the local eligibility surface.  Combining with $|\Geff|\sim m_0^2M_G^2/\alpha_G$ gives
\begin{equation}
\frac{m_0^2}{\alpha_G M_G^2}\ll1.
\label{eq:EFTcontrol}
\end{equation}
Thus activation can occur inside the EFT regime if $m_0\ll \sqrt{\alpha_G}M_G$.


\subsection{Linear spectral test: local negativity is not enough}
To determine whether the symmetric state $H=0$ actually destabilizes, set $A_\mu=0$ and linearize $H=\epsilon\,\delta H+O(\epsilon^2)$.  In Schwarzschild, decompose
\begin{equation}
\delta H=e^{-i\omega t}Y_{\ell m}(\theta,\phi)\frac{u_\ell(r)}{r}.
\end{equation}
With $f(r)=1-2M/r$ and $\dd r_*/\dd r=f^{-1}$, the exact linear equation becomes
\begin{equation}
\left[-\frac{\dd^2}{\dd r_*^2}+V_\ell(r)\right]u_\ell=\omega^2 u_\ell,
\label{eq:spectral}
\end{equation}
where
\begin{equation}
\boxed{
V_\ell=f\left[
\frac{\ell(\ell+1)}{r^2}+\frac{2M}{r^3}+m_0^2
-\frac{48\alpha_GM^2}{M_G^2r^6}
\right].}
\label{eq:Vschw}
\end{equation}
A tachyonic mode exists when the lowest eigenvalue has $\omega^2<0$.  Thus Eq.~\eqref{eq:rG} is only a local eligibility radius; the spectrum of Eq.~\eqref{eq:spectral} is the correct onset criterion.

The distinction admits a useful variational test.  Define
\begin{equation}
x=\frac{r}{2M},\qquad
\mu=2Mm_0,\qquad
\gamma=\frac{3\alpha_G}{M_G^2M^2},\qquad f=1-\frac1x.
\end{equation}
Then $(2M)^2V_\ell=f B_\ell$ with
\begin{equation}
B_\ell(x)=\frac{\ell(\ell+1)}{x^2}+\frac1{x^3}+\mu^2-\frac{\gamma}{x^6}.
\end{equation}
For any admissible trial function $u$, the Rayleigh quotient is
\begin{equation}
\boxed{
\mathcal R[u]=
\frac{\displaystyle\int_1^\infty\!\dd x\,
\left[f|u'|^2+B_\ell|u|^2\right]}
{\displaystyle\int_1^\infty\!\dd x\,|u|^2/f}.}
\label{eq:rayleigh}
\end{equation}
If $\mathcal R[u]<0$ for even one trial function, the lowest mode is necessarily unstable.  The accompanying check script evaluates Eq.~\eqref{eq:rayleigh}; for example, for $\ell=0$, $\mu=0.2$, $\gamma=12$, and $u=(x-1)e^{-2(x-1)}$, it finds $\mathcal R\simeq-6.89\times10^{-2}$.  This does not locate the exact critical coupling, but it demonstrates a parameter point where the model passes a genuine global instability test rather than only the local sign test.

\section{Rotational curvature: a covariant ``tornado'' variable}
Spin-triggered black-hole scalarization is already established in Gauss--Bonnet models, and quadratic couplings to the Pontryagin density have also been studied \citep{HerdeiroEtAl2021,DonevaCollodelYazadjiev2022,GaoHuangLiu2019}.  Rotation is therefore not a novelty claim.  Its role here is to sharpen the weather-inspired trigger while retaining covariance and parity symmetry.

For Kerr in Boyer--Lindquist coordinates, define
\begin{equation}
\rho^2=r^2+a^2\cos^2\theta.
\end{equation}
Because Kerr is Ricci flat, the Gauss--Bonnet density equals the Kretschmann scalar,
\begin{equation}
\Geff_{\rm K}=
\frac{48M^2\left(r^6-15a^2r^4c^2+15a^4r^2c^4-a^6c^6\right)}{\rho^{12}},
\qquad c\equiv\cos\theta,
\label{eq:GKerr}
\end{equation}
and the Pontryagin density is, up to the orientation convention,
\begin{equation}
\Pontryagin_{\rm K}=
-\frac{96aM^2rc\left(r^2-3a^2c^2\right)\left(3r^2-a^2c^2\right)}{\rho^{12}}.
\label{eq:PKerr}
\end{equation}
Equation~\eqref{eq:PKerr} is odd across the equatorial plane and vanishes for $a=0$; its square is parity even.  These Kerr invariants and Pontryagin-based scalar couplings are standard \citep{GaoHuangLiu2019,ZhouChenJing2021}; the new question is what they do when embedded in the finite-state defect sector.

Ignoring the dark term, the local weather surface is defined by
\begin{equation}
m_0^2-\frac{\alpha_G}{M_G^2}\Geff_{\rm K}
-\frac{\alpha_\Omega}{M_\Omega^6}\Pontryagin_{\rm K}^2=0.
\label{eq:kerrweather}
\end{equation}
For slow rotation,
\begin{align}
\Geff_{\rm K}&=\frac{48M^2}{r^6}
\left[1-21\frac{a^2c^2}{r^2}+O(a^4/r^4)\right],\\
\Pontryagin_{\rm K}^2&=\frac{82944M^4a^2c^2}{r^{14}}+O(a^4/r^{16}).
\end{align}
Let
\begin{equation}
r_0=\left(\frac{48\alpha_GM^2}{m_0^2M_G^2}\right)^{1/6}
\end{equation}
be the Schwarzschild local radius.  Solving Eq.~\eqref{eq:kerrweather} perturbatively gives
\begin{equation}
\boxed{
\frac{r_{\rm w}(\theta)}{r_0}
=1+\frac{a^2\cos^2\theta}{r_0^2}
\left(-\frac72+6\eta_\Omega\right)
+O(a^4/r_0^4),}
\label{eq:weatherdeform}
\end{equation}
where
\begin{equation}
\boxed{
\eta_\Omega\equiv
\frac{\alpha_\Omega m_0^2M_G^4}{\alpha_G^2M_\Omega^6}.}
\label{eq:etaOmega}
\end{equation}
The black-hole mass cancels from $\eta_\Omega$.  Consequently the sign of the leading polar deformation is controlled by a universal EFT ratio:
\begin{equation}
\eta_\Omega>\frac7{12}
\quad\Longrightarrow\quad
r_{\rm w}(0)>r_0
\quad\text{at leading order},
\label{eq:shapecriterion}
\end{equation}
while $\eta_\Omega<7/12$ gives the opposite polar shift.  This is a concrete weather-inspired prediction: rotation does not merely ``look like'' a tornado; it deforms an invariant phase-eligibility surface in a calculable way.  Equation~\eqref{eq:weatherdeform} remains a local criterion.  The actual Kerr onset requires solving the two-dimensional nonseparable spectral problem associated with Eq.~\eqref{eq:kerrweather}.

\section{Dark-halo eligibility from observable density profiles}

The dark-selective term implements the motivating possibility that hidden mass can shift where defects become dynamically possible.  In a weak-curvature halo, Eq.~\eqref{eq:activation} reduces to
\begin{equation}
\Ddark(r)>\rho_c,
\qquad
\boxed{
\rho_c\equiv\frac{m_0^2M_D^2}{\alpha_D}.
}
\label{eq:rhoc}
\end{equation}
For nonrelativistic dark matter $\Ddark\simeq\rho_D$.

Take an NFW profile
\begin{equation}
\rho_D(r)=\frac{\rho_s}{x(1+x)^2},
\qquad x\equiv\frac{r}{r_s}.
\end{equation}
The local eligibility radius $r_c=r_sx_c$ satisfies
\begin{equation}
\frac{\rho_s}{x_c(1+x_c)^2}=\rho_c,
\end{equation}
so
\begin{equation}
\boxed{
x_c(1+x_c)^2=\frac{\rho_s}{\rho_c}.}
\label{eq:nfwthreshold}
\end{equation}
For each halo, $(\rho_s,r_s)$ may be inferred from dynamical or lensing analyses.  A universal parameter $\rho_c$ then predicts the radius inside which the local potential favors the discrete Higgs phase.  The model is falsified if a single universal parameter set cannot consistently describe the chosen class of systems once a concrete observable of the defect sector is specified.

Equation~\eqref{eq:nfwthreshold} predicts an \emph{eligible region}, not a deterministic number of strings.  Defect abundance depends on phase-transition history, correlation length, and annihilation.  Claiming otherwise would overstate the theory.

\section{What is actually different from nearby literature?}

The individual ingredients have substantial precedent:
\begin{itemize}[leftmargin=1.5em]
    \item Einstein dynamics has a well-known thermodynamic interpretation based on local horizon entropy \citep{Jacobson1995,ElingGuedensJacobson2006}.
    \item The membrane paradigm gives black-hole horizons fluid-like stress and transport properties \citep{ParikhWilczek1998}.
    \item Density-dependent scalar symmetry breaking and screening are central to symmetron models \citep{HinterbichlerKhoury2010}.
    \item Matter-dependent local $U(1)$ cosmic strings already occur in the complex symmetron model, and simulations show halo attachment \citep{NezhadsafaviPogosian2025}.
    \item Environment-dependent string and wall networks have been evolved in relativistic cosmological simulations \citep{ChristiansenAdamekKunz2026}.
    \item Quantized vortices in bosonic dark matter are established \citep{KainLing2010,BraxValageas2025}.
    \item Curvature- and spin-triggered scalar phases around black holes are established, including Kerr scalarization in Gauss--Bonnet models \citep{DonevaYazadjiev2018,HerdeiroEtAl2021,DonevaCollodelYazadjiev2022,LaiEtAl2026}.
    \item Quadratic scalar couplings to the Kerr Pontryagin density can also generate tachyonic instabilities \citep{GaoHuangLiu2019}.
    \item Discrete black-hole spectra and degeneracy arguments have a long history \citep{BekensteinMukhanov1995}.
    \item $\mathbb Z_p$ gauge theories and BF descriptions are standard \citep{KapustinSeiberg2014,GaiottoEtAl2015}.
\end{itemize}

The proposed novelty is therefore deliberately narrower.  It is the conjunction of the following statements in one covariant EFT:
\begin{enumerate}[leftmargin=1.7em]
    \item high dark-sector trace and/or curvature can \emph{induce} a hidden Higgs phase rather than restore it;
    \item rotational curvature enters through a parity-even $\Pontryagin^2$ trigger, producing the analytic slow-spin deformation in Eq.~\eqref{eq:weatherdeform};
    \item the Higgs has charge $p$, leaving a local residual $\mathbb Z_p$ phase rather than a continuous long-range Goldstone sector;
    \item the same integer $p$ appears in an exact macroscopic holonomy, Eq.~\eqref{eq:AB}; and
    \item the onset question is separated cleanly into local eligibility surfaces and a global spectral instability, Eq.~\eqref{eq:spectral}, so the model is not allowed to identify $m_E^2<0$ with defect formation by fiat.
\end{enumerate}

We have not found, in the literature audited for this draft, a peer-reviewed paper with precisely this combined activation law and residual discrete gauge observable.  This is a \emph{novelty hypothesis}, not proof of priority; a broader dedicated literature search would still be required before submission.

\section{Thermodynamic consistency target}

Calling the sector ``entropic'' should ultimately be justified by state counting.  A minimal consistency check is the following.  Suppose a horizon patch of area $A$ contains $N=A/a_\ast$ independent microscopic cells, each with $p$ accessible states.  Then
\begin{equation}
\Omega=p^N,
\qquad
S=k_B\ln\Omega
=k_B\frac{A}{a_\ast}\ln p.
\end{equation}
Matching the Bekenstein--Hawking law
\begin{equation}
S_{\rm BH}=k_B\frac{A}{4\ell_P^2}
\end{equation}
requires
\begin{equation}
\boxed{
a_\ast=4\ell_P^2\ln p.}
\label{eq:areacell}
\end{equation}
Equation~\eqref{eq:areacell} is only a matching condition.  Discrete black-hole area and degeneracy constructions predate this work \citep{BekensteinMukhanov1995}; we do not claim Eq.~\eqref{eq:areacell} as new.  Its value here is diagnostic: a microscopic completion of the $\mathbb Z_p$ sector should explain, not merely fit, any horizon entropy relation.

\section{Hard audit: failure modes}

The model should be abandoned or substantially revised if any of the following occurs.

\begin{enumerate}[leftmargin=1.7em]
    \item \textbf{Literature collision.}  An existing model is found with the same high-density/curvature activation, charge-$p$ Higgsing, and the same proposed observational channel.
    \item \textbf{No controlled EFT window.}  Eligibility requires $|\Geff|/M_G^4\gtrsim1$, $|\Pontryagin|/M_\Omega^4\gtrsim1$, or strong coupling before the defect phase is reached.  In particular the $\Pontryagin^2$ operator is high dimension and must not be treated nonperturbatively beyond its cutoff.
    \item \textbf{No global instability.}  If the spectral problem has no unstable mode anywhere in the EFT-controlled region, the local negative-mass surfaces are physically insufficient.
    \item \textbf{No viable cosmology.}  Defect energy density, gravitational radiation, or dark-sector scattering exceeds established limits for every parameter region that gives astrophysically relevant activation.
    \item \textbf{No entropy derivation.}  The $\mathbb Z_p$ degree of freedom remains a generic hidden gauge sector with no independent reason to call it entropic.  In that case the model may still be a dark-sector defect EFT, but the original entropy interpretation should be removed.
    \item \textbf{No distinctive observable.}  If all effects reduce to an ordinary Nielsen--Olesen string with redefined parameters and the phase in Eq.~\eqref{eq:AB} is inaccessible, the finite arithmetic has no empirical leverage.
    \item \textbf{Pressure intuition fails.}  If solutions of Eq.~\eqref{eq:pressure} do not exhibit the proposed pressure depression, the ``storm eye'' language must be dropped rather than tuned into existence.
\end{enumerate}

\section{Concrete next derivations}

The shortest path from framework to a serious paper is now sharply defined:
\begin{enumerate}[leftmargin=1.7em]
    \item Solve Eq.~\eqref{eq:spectral} for the Schwarzschild critical curve $\gamma_c(\mu)$ and then solve the full Kerr $(r,\theta)$ eigenvalue problem.  The slow-spin expression~\eqref{eq:weatherdeform} should emerge as the local geometric limit, not replace the spectrum.
    \item Solve the nonlinear Abelian-Higgs equations in an inhomogeneous background and compute $\rho_E$, radial pressure, azimuthal pressure, string tension, and whether a genuine low-pressure core exists.
    \item Determine whether the $\mathbb Z_p$ holonomy produces a measurable $p$-dependent Aharonov--Bohm scattering or transport effect for a specified hidden-matter probe.  Without this step, finite arithmetic remains formal.
    \item Combine lensing-inferred halo profiles with the global field equation rather than only Eq.~\eqref{eq:nfwthreshold}, testing whether one universal parameter set predicts where the hidden phase can actually condense.
    \item Perform a parameter-space audit including fifth-force avoidance, defect abundance, dark self-interaction bounds, EFT control, black-hole horizon accessibility, and existing scalarization/ringdown constraints.
\end{enumerate}

A journal-level result would require at least the first two calculations plus one distinctive observable.  An arXiv research note could reasonably precede the full nonlinear program only if the spectral critical curve and the Kerr deformation survive independent checking.

\section{Conclusion}

The weather metaphor is retained only where it generates invariant structure.  Circular matter flow can be vortex-like, but an orbit is not itself a spacetime tornado.  Rotational spacetime structure is instead encoded through Kerr curvature invariants.  The resulting compact EFT contains a charge-$p$ hidden Higgs field whose phase is influenced by Gauss--Bonnet curvature, parity-even rotational curvature $\Pontryagin^2$, and a dark-selective stress-tensor trace.  Its finite content survives as the exact $\mathbb Z_p$ holonomy $\exp(2\pi i n/p)$.

Two substantive corrections distinguish this version from a metaphor-driven model.  First, $m_E^2<0$ is only a local eligibility condition; the Schwarzschild onset is governed by the spectral operator in Eq.~\eqref{eq:spectral}, with Eq.~\eqref{eq:rayleigh} providing a rigorous sufficient test for instability.  Second, Kerr rotation produces the explicit slow-spin deformation~\eqref{eq:weatherdeform}, whose leading polar sign flips at the mass-independent coupling ratio $\eta_\Omega=7/12$.  These are calculable statements, not analogical identifications.

The project is still not a completed theory of spacetime entropy.  Its strongest present status is a falsifiable hidden-sector EFT research program with a narrow novelty hypothesis.  The decisive next step is to compute the true Schwarzschild/Kerr instability boundaries and nonlinear defect solutions; if those do not yield a controlled parameter window and a $p$-dependent observable, the ``entropic weather'' interpretation should be abandoned.

\printbibliography

\end{document}
